\documentclass[11pt]{article}
\usepackage[spanish,es-noshorthands]{babel}
\usepackage[utf8]{inputenc}
\usepackage[T1]{fontenc}
\usepackage{amsmath,amssymb,amsthm,mathtools,booktabs,geometry,enumitem,xcolor}
\usepackage{algorithm}
\usepackage{algpseudocode}
\usepackage[hidelinks]{hyperref}
\geometry{margin=2.4cm}
\numberwithin{equation}{section}
\newtheorem{proposition}{Proposición}
\theoremstyle{remark}
\newtheorem*{remark}{Observación}
\newcommand{\bu}{\mathbf{u}}
\newcommand{\bK}{\mathbf{K}}
\newcommand{\bD}{\mathbf{D}}
\newcommand{\bn}{\mathbf{n}}
\newcommand{\bb}{\mathbf{b}}
\newcommand{\bd}{\mathbf{d}}
\newcommand{\bx}{\mathbf{x}}
\newcommand{\Ceps}{g_\varepsilon}
\newcommand{\im}{\mathrm{im}}
\newcommand{\pA}{\psi^{\mathrm A}}          % multiplicador ADR
\newcommand{\pM}{\psi^{\mathrm M}}          % multiplicador MRMT (móvil)
\newcommand{\pim}{\psi_{\im}}               % multiplicador región inmóvil
\newcommand{\phA}{\psi_h^{\mathrm A}}
\newcommand{\phM}{\psi_h^{\mathrm M}}
\definecolor{cSum}{RGB}{200,30,30}   % sumidero  fC
\definecolor{cDiv}{RGB}{25,80,200}   % divergencia  C div u
\definecolor{cAlm}{RGB}{20,140,60}   % almacenamiento  S_s h_t C
\newcommand{\sumi}[1]{{\color{cSum}#1}}
\newcommand{\divi}[1]{{\color{cDiv}#1}}
\newcommand{\alm}[1]{{\color{cAlm}#1}}
\algrenewcommand\algorithmicrequire{\textbf{Entrada:}}
\algrenewcommand\algorithmicensure{\textbf{Salida:}}
\algrenewcommand\algorithmicfor{\textbf{para}}
\algrenewcommand\algorithmicdo{\textbf{hacer}}
\algrenewcommand\algorithmicwhile{\textbf{mientras}}
\algrenewcommand\algorithmicif{\textbf{si}}
\algrenewcommand\algorithmicthen{\textbf{entonces}}
\algrenewcommand\algorithmicreturn{\textbf{devolver}}
\algrenewcommand\algorithmicend{\textbf{fin}}
\floatname{algorithm}{Algoritmo}

\title{Modelo ADR/MRMT sin sumidero de contaminante:\\
eliminación de términos, adjuntos por variante y\\ algoritmos de gradiente conjugado}
\author{Tesis doctoral --- control óptimo de la contaminación en acuíferos}
\date{}

\begin{document}
\maketitle

\begin{abstract}
Se parte del balance de soluto en forma conservativa, conservando
$\nabla\!\cdot\mathbf{u}$, y se demuestra que el sumidero explícito
$\chi(Q)g_\varepsilon C$ queda cancelado y que sobrevive el término de
almacenamiento $-S_s\partial_th\,C$. De ahí salen las ecuaciones del modelo
nuevo para ADR y para MRMT. Para cada variante se construye el Lagrangiano,
se deduce el sistema adjunto (transporte, regiones inmóviles y flujo) y el
gradiente respecto de $Q(t)$ y $z_p$; se escribe el adjunto discreto de ADR,
MRMT por bloque y MRMT semi-analítico, y los algoritmos de gradiente conjugado
que los usan.
\end{abstract}

%===========================================================
\section{Punto de partida}
\label{sec:partida}
%===========================================================

\paragraph{Código de colores.}
$\sumi{\text{rojo}}$: sumidero $fC$ (el que se elimina);
$\divi{\text{azul}}$: divergencia $C\,\nabla\!\cdot\bu$ (de donde sale el sumidero);
$\alm{\text{verde}}$: almacenamiento $S_s\,\partial_th\,C$ (el que queda).

\paragraph{Hoja de ruta.} (i) §1: ecuaciones de flujo y transporte conservativo. (ii) §2: se desarrolla
$\nabla\!\cdot(\bu C)$, se sustituye $\nabla\!\cdot\bu$ y se cancela el sumidero; salen N-ADR y N-MRMT.
(iii) §3: se define $J$ y un Lagrangiano por variante. (iv) §4--5: se deriva el Lagrangiano respecto de
$C$, de $C_{\im,r}$ y de $h$ (cada derivada, término a término) y se anula; salen los adjuntos y el gradiente.
(v) §6: lo mismo en forma discreta. (vi) §7: los algoritmos.

\paragraph{Notación.}
\begin{center}\small
\begin{tabular}{p{3.3cm}p{11.2cm}}
\toprule
símbolo & significado \\
\midrule
$\Omega,\ \partial\Omega,\ \bn$ & dominio $\subset\mathbb{R}^2$, su frontera y la normal exterior; $Q_T=\Omega\times(0,T)$ \\
$\Gamma_{\rm in},\Gamma_{\rm out},\Gamma_{\rm wall}$ & entrada (derecha), salida (izquierda) y paredes; $\partial\Omega$ es su unión \\
$h,\ \bu=-\bK\nabla h$ & carga hidráulica y velocidad de Darcy ($\bK$: conductividad) \\
$S_s,\ \phi,\ R$ & almacenamiento específico, porosidad y retardo (constantes en el tiempo) \\
$C$ ($C_m$ en MRMT), $C_{\im,r}$ & concentración móvil y en la región inmóvil $r=1,\dots,N_r$ \\
$\bD(\bu)$ & dispersión, $\bD=\bD(\bu)$ (depende de la velocidad) \\
$\lambda_d,\ \lambda_r,\ \phi_{\im,r}$ & decaimiento, tasa de intercambio y porosidad inmóvil de la región $r$ \\
$c_0$ & coeficiente de reacción: $\lambda_d\phi R$ (ADR), $\sum_r\lambda_r$ (MRMT) \\
$Q(t),\ z_p$ & caudal de bombeo y profundidad del pozo (variables de decisión) \\
$g_\varepsilon(\cdot;z_p)$ & kernel gaussiano de masa unidad centrado en el pozo; $\dot g_\varepsilon=\partial_{z_p}g_\varepsilon$ \\
$f=-Q\,g_\varepsilon$ & sumidero de flujo (negativo: extracción); $f_{\rm in}$: flujo prescrito en $\Gamma_{\rm in}$ \\
$W_+$ & fuente de contaminante \\
$(\cdot,\cdot)=\langle\cdot,\cdot\rangle$ & producto escalar de $L^2(\Omega)$; $\oint=\int_{\partial\Omega}$ \\
$\psi^{\rm A},\psi^{\rm M},\psi_{\im}^{\,r}$ & multiplicadores (adjuntos) del transporte ADR, del transporte móvil MRMT y de la región inmóvil $r$ \\
$\psi_h^{\rm A},\psi_h^{\rm M}$ & multiplicadores (adjuntos) del flujo en cada variante \\
$\hat J(Q,z_p)$ & funcional reducido (se evalúa resolviendo el estado) \\
\bottomrule
\end{tabular}
\end{center}
Convención de colores y de subíndices: $\psi^{\rm A}$ y $\psi^{\rm M}$ son el multiplicador de la ecuación de $C$
(el superíndice indica la variante); el subíndice $h$ indica el multiplicador de la ecuación de $h$.

\paragraph{Flujo.}
\begin{equation}
S_s\,\partial_th-\nabla\!\cdot(\bK\nabla h)=f=-Q\,\Ceps,\qquad \bu=-\bK\nabla h .
\label{eq:flujo}
\end{equation}
Como $\nabla\!\cdot\bu=-\nabla\!\cdot(\bK\nabla h)$, la divergencia se despeja del flujo:
\begin{equation}
\boxed{\;\divi{\nabla\!\cdot\bu}\;=\;\sumi{f}\;-\;\alm{S_s\,\partial_th}
\;=\;\sumi{-Q\,\Ceps}\;-\;\alm{S_s\,\partial_th}\;}
\label{eq:div_u}
\end{equation}

\paragraph{Transporte en forma conservativa.}
La masa por unidad de volumen es $\phi RC$, el flujo de soluto es $\bu C-\bD\nabla C$,
y el agua que extrae el pozo se lleva el soluto con la concentración local
(retira masa a razón $fC$):
\begin{equation}
\phi R\,\partial_tC+\nabla\!\cdot(\bu\,C-\bD\nabla C)+\lambda_d\phi R\,C
=W_++\sumi{f\,C}.
\label{eq:conservativa}
\end{equation}

%===========================================================
\section{Eliminación de términos y ecuaciones del modelo nuevo}
\label{sec:eliminacion}
%===========================================================

\subsection{ADR}
\paragraph{Paso 1: desarrollar sin descartar $\nabla\!\cdot\bu$.}
$\nabla\!\cdot(\bu C)=\bu\cdot\nabla C+\divi{C\,\nabla\!\cdot\bu}$, luego
\begin{equation}
\phi R\,\partial_tC+\bu\cdot\nabla C-\nabla\!\cdot(\bD\nabla C)+\lambda_d\phi RC
+\divi{C\,\nabla\!\cdot\bu}=W_++\sumi{f\,C}.
\label{eq:desarrollada}
\end{equation}

\paragraph{Paso 2: sustituir la divergencia despejada del flujo.}
Por \eqref{eq:div_u}, $\divi{C\nabla\!\cdot\bu}=\sumi{fC}-\alm{S_s\partial_th\,C}$:
\begin{equation}
\phi R\,\partial_tC+\bu\cdot\nabla C-\nabla\!\cdot(\bD\nabla C)+\lambda_d\phi RC
+\sumi{fC}-\alm{S_s\partial_th\,C}=W_++\sumi{fC}.
\end{equation}

\paragraph{Paso 3: cancelación.} El $\sumi{fC}$ está en ambos miembros y se anula;
el almacenamiento no:
\begin{equation}
\boxed{\;\phi R\,\partial_tC+\bu\cdot\nabla C-\nabla\!\cdot(\bD\nabla C)
+\lambda_d\phi R\,C\;\alm{-\,S_s\,\partial_th\;C}=W_+\;}
\tag{N-ADR}
\label{eq:N_ADR}
\end{equation}

\begin{center}\small
\begin{tabular}{p{2.6cm}p{5.4cm}p{6.6cm}}
\toprule
término & origen & destino \\
\midrule
$\sumi{fC}$ & miembro derecho de \eqref{eq:conservativa} (extracción) & se cancela con la parte $\sumi{fC}$ de $\divi{C\nabla\!\cdot\bu}$ \\
$\divi{C\nabla\!\cdot\bu}$ & regla del producto en $\nabla\!\cdot(\bu C)$ & $=\sumi{fC}-\alm{S_s\partial_th\,C}$ por \eqref{eq:div_u} \\
$\alm{S_s\partial_th\,C}$ & parte de $\nabla\!\cdot\bu$ debida al almacenamiento & \textbf{queda} en \eqref{eq:N_ADR} \\
\bottomrule
\end{tabular}
\end{center}

\begin{proposition}[Equivalencia]\label{prop:equiv}
Sea $h$ solución de \eqref{eq:flujo}. $C$ resuelve \eqref{eq:conservativa} si y solo si
resuelve \eqref{eq:N_ADR}.
\end{proposition}
\begin{proof}
Los pasos 1--3 son identidades reversibles y del flujo solo se usa \eqref{eq:div_u}.
En sentido débil la cuenta es la misma probando con $v\in H^1(\Omega)$ (no hay término de
frontera en la forma no integrada por partes de $\bu\cdot\nabla C$).
\end{proof}

\subsection{MRMT}
El intercambio con las regiones inmóviles, $\sum_r\lambda_r(C_m-C_{\im,r})$, no contiene
$\bu$ ni $f$, y la ecuación inmóvil es una EDO local ($\phi_{\im,r}$ constante). Los pasos 1--3
se repiten sin cambios sobre la ecuación móvil
$\phi R\,\partial_tC_m+\nabla\!\cdot(\bu C_m-\bD\nabla C_m)+\sum_r\lambda_r(C_m-C_{\im,r})=W_++\sumi{fC_m}$:
\begin{equation}
\boxed{\begin{aligned}
&\phi R\,\partial_tC_m+\bu\cdot\nabla C_m-\nabla\!\cdot(\bD\nabla C_m)
+\sum_{r=1}^{N_r}\lambda_r\,(C_m-C_{\im,r})\;\alm{-\,S_s\partial_th\,C_m}=W_+,\\
&\phi_{\im,r}R\,\partial_tC_{\im,r}=\lambda_r\,(C_m-C_{\im,r}),\quad r=1,\dots,N_r.
\end{aligned}}
\tag{N-MRMT}
\label{eq:N_MRMT}
\end{equation}
En lo que sigue se escribe $C$ por $C_m$ en MRMT y se define el coeficiente de reacción
$c_0=\lambda_d\phi R$ en ADR y $c_0=\sum_r\lambda_r$ en MRMT.

\subsection{Condiciones de frontera e iniciales}
Con $\Gamma_{\rm in}$ a la derecha, $\Gamma_{\rm out}$ a la izquierda y $\Gamma_{\rm wall}$ el
resto (el flujo va de derecha a izquierda):
\begin{center}\small
\begin{tabular}{lll}
\toprule
 & flujo $h$ & transporte $C$ \\
\midrule
$\Gamma_{\rm in}$ & $(\bK\nabla h)\cdot\bn=f_{\rm in}\ge0$, es decir $\bu\cdot\bn=-f_{\rm in}\le0$ & $(\bD\nabla C)\cdot\bn=0$ \\
$\Gamma_{\rm wall}$ & $(\bK\nabla h)\cdot\bn=0$ & $(\bD\nabla C)\cdot\bn=0$ \\
$\Gamma_{\rm out}$ & libre ($\nabla(\bK\nabla h)\cdot\bn=0$) & $(\bD\nabla C)\cdot\bn=0$ \\
\bottomrule
\end{tabular}
\end{center}
El contaminante solo cruza la frontera por advección. Condiciones iniciales:
$h(0)=h_0$, $C(0)=C_0$, $C_{\im,r}(0)=C_{\im,r,0}$.

%===========================================================
\section{Funcional y Lagrangianos}
\label{sec:lagrangiano}
%===========================================================

\begin{equation}
J=\gamma J_m+\alpha J_{\rm inst}+\beta J_{\rm op},\qquad
J_m=\!\int_0^T\!\!\!\int_\Omega\Ceps\,C^2,\quad
J_{\rm inst}=\kappa_p(z_p-z_0)^2,\quad J_{\rm op}=d_p\!\int_0^TQ^2\,dt,
\end{equation}
con $d_p=z_0-z_p$. Espacios $V_h=V_C=H^1(\Omega)$, $V_{\im}=L^2(\Omega)$. Formas del flujo:
\begin{equation}
a_h(h,v)=\int_\Omega\bK\nabla h\cdot\nabla v-\int_{\Gamma_{\rm out}}(\bK\nabla h\cdot\bn)v,\qquad
\ell_h(v;Q,z_p)=-Q\!\int_\Omega\Ceps v+\int_{\Gamma_{\rm in}}f_{\rm in}v .
\end{equation}
Forma del transporte (común a ambas variantes salvo $c_0$ y el acoplamiento):
\begin{equation}
a_C(C,v;h)=\int_\Omega\bD(\bu)\nabla C\cdot\nabla v+\int_\Omega(\bu\cdot\nabla C)\,v
+\int_\Omega c_0\,Cv\;\alm{-\int_\Omega S_s\,\partial_th\;C\,v}.
\label{eq:aC_nueva}
\end{equation}
$a_C$ no depende del control, pero sí de $h$, por $\bu$ y por $\partial_th$.

\paragraph{Estado en forma débil (lo que se multiplica por los multiplicadores).} Para todo $v\in V_h$ y
todo $w\in V_C$, en cada instante:
\begin{equation}
\underbrace{\langle S_s\partial_th,v\rangle+a_h(h,v)-\ell_h(v;Q,z_p)}_{R^h(h;v)}=0,\qquad
\underbrace{\langle\phi R\partial_tC,w\rangle+a_C(C,w;h)-(W_+,w)}_{R^C(C,h;w)}=0 .
\label{eq:estado_debil}
\end{equation}
\emph{Idea del método.} Se resuelve el problema $\min\hat J$ sujeto a $R^h=0$ y $R^C=0$ añadiendo a $J$ los
residuos multiplicados por $\psi_h$ y $\psi_C$ (Lagrangiano). Si $h$ y $C$ cumplen el estado, el
Lagrangiano vale $J$; si además se eligen $\psi_h,\psi_C$ de modo que $\partial_C\mathcal{L}=0$ y
$\partial_h\mathcal{L}=0$, la derivada de $\hat J$ respecto del control es la derivada parcial de
$\mathcal{L}$ respecto del control, sin tener que derivar el estado.
Convención: los residuos se suman a $J$.

\paragraph{Lagrangiano ADR.}
\begin{multline}
\mathcal{L}^{\rm A}=J+\!\int_0^T\!\!\Bigl[\langle S_s\partial_th,\phA\rangle+a_h(h,\phA)-\ell_h(\phA;Q,z_p)\\
+\langle\phi R\partial_tC,\pA\rangle+a_C(C,\pA;h)-(W_+,\pA)\Bigr]dt .
\label{eq:L_ADR}
\end{multline}

\paragraph{Lagrangiano MRMT.}
\begin{multline}
\mathcal{L}^{\rm M}=J+\!\int_0^T\!\!\Bigl[\langle S_s\partial_th,\phM\rangle+a_h(h,\phM)-\ell_h(\phM;Q,z_p)
+\langle\phi R\partial_tC_m,\pM\rangle+a_C(C_m,\pM;h)\\
-\textstyle\sum_r\lambda_r(C_{\im,r},\pM)-(W_+,\pM)
+\sum_r\bigl(\phi_{\im,r}R\partial_tC_{\im,r}+\lambda_r(C_{\im,r}-C_m),\pim^{\,r}\bigr)\Bigr]dt .
\label{eq:L_MRMT}
\end{multline}
Las variaciones respecto de los multiplicadores recuperan el estado. Como $J$ no tiene término
terminal, $\psi(T)=0$ en todos los multiplicadores.

%===========================================================
\section{Adjunto del modelo ADR}
\label{sec:adj_ADR}
%===========================================================

\subsection{Variación respecto de $C$: adjunto del transporte}
Se perturba $C\to C+\epsilon\,\eta$ con $\eta\in V_C$, $\eta(0)=0$ (el estado inicial es dato), y se
calcula $\frac{d}{d\epsilon}\mathcal{L}^{\rm A}\big|_{\epsilon=0}$ término a término:
\begin{enumerate}[label=\textbf{(\alph*)},nosep]
\item \emph{Funcional.} $J_m=\int\!\!\int\Ceps C^2$ aporta $2\gamma\int_0^T(\Ceps C,\eta)\,dt$ (los términos
  $J_{\rm inst}$, $J_{\rm op}$ y $R^h$ no dependen de $C$).
\item \emph{Derivada temporal.} $\int_0^T\langle\phi R\partial_t\eta,\pA\rangle dt
  =\bigl[\langle\phi R\eta,\pA\rangle\bigr]_0^T-\int_0^T\langle\phi R\,\eta,\partial_t\pA\rangle dt
  =-\int_0^T\langle\phi R\,\partial_t\pA,\eta\rangle dt$, porque $\eta(0)=0$ y $\pA(T)=0$.
\item \emph{Forma $a_C$.} Es lineal en $C$, así que aporta $\int_0^Ta_C(\eta,\pA;h)\,dt$.
\item \emph{Fuente $W_+$.} No depende de $C$ y no aporta.
\end{enumerate}
Sumando (a)--(c) e imponiendo que sea cero para toda $\eta$ (y todo instante) se obtiene el problema débil
\begin{equation}
-\langle\phi R\partial_t\pA,\eta\rangle+a_C(\eta,\pA;h)=-2\gamma(\Ceps C,\eta)\quad\forall\eta\in V_C,
\qquad\pA(T)=0 .
\label{eq:adj_debil}
\end{equation}
\paragraph{Paso a forma fuerte.} Se integra por partes en espacio los dos términos de $a_C(\eta,\pA;h)$ que
tienen derivadas de $\eta$ (los de orden cero, $\int c_0\eta\pA$ y $-\int S_s\partial_th\,\eta\pA$, pasan tal cual):
\[
\int_\Omega\bD\nabla\eta\cdot\nabla\pA=-\!\int_\Omega\eta\,\nabla\!\cdot(\bD\nabla\pA)+\oint\eta\,(\bD\nabla\pA)\cdot\bn,
\qquad
\int_\Omega(\bu\cdot\nabla\eta)\pA=-\!\int_\Omega\eta\,\nabla\!\cdot(\bu\pA)+\oint(\bu\cdot\bn)\eta\,\pA .
\]
Se separa en dos pasos: (1) con $\eta$ de soporte compacto en $\Omega$ los términos de borde se anulan
y queda la ecuación en el interior; (2) con $\eta$ arbitraria en $\partial\Omega$ (no hay Dirichlet),
la suma de los dos términos de borde, $\oint\eta\,(\bD\nabla\pA+\bu\,\pA)\cdot\bn$, debe ser cero:
es la condición natural del adjunto (flujo total nulo de $\pA$). Resultado:
\begin{equation}
\boxed{\begin{aligned}
&-\phi R\,\partial_t\pA-\nabla\!\cdot(\bD\nabla\pA)-\nabla\!\cdot(\bu\,\pA)
+c_0\,\pA\;\alm{-\,S_s\partial_th\,\pA}=-2\gamma\,\Ceps\,C,\\
&\pA(\cdot,T)=0,\qquad(\bD\nabla\pA+\bu\,\pA)\cdot\bn=0\ \text{ en }\partial\Omega .
\end{aligned}}
\label{eq:adj_fuerte}
\end{equation}
En $\Gamma_{\rm wall}$, como $\bu\cdot\bn=0$, la condición se reduce a $(\bD\nabla\pA)\cdot\bn=0$.
Es el adjunto del flujo dispersivo nulo del estado.

\subsection{Eliminación en el adjunto: forma no conservativa}
$-\nabla\!\cdot(\bu\pA)=-\bu\cdot\nabla\pA-\divi{\pA\,\nabla\!\cdot\bu}$ y, por \eqref{eq:div_u},
$-\divi{\pA\nabla\!\cdot\bu}=\sumi{Q\Ceps\pA}+\alm{S_s\partial_th\,\pA}$. Sustituyendo:
\begin{equation}
-\phi R\,\partial_t\pA-\nabla\!\cdot(\bD\nabla\pA)-\bu\cdot\nabla\pA
+c_0\pA+\sumi{Q\,\Ceps\,\pA}
\;\underbrace{\alm{+\,S_s\partial_th\,\pA-S_s\partial_th\,\pA}}_{=\,0}=-2\gamma\,\Ceps\,C .
\label{eq:adj_noconservativa}
\end{equation}
El verde de la divergencia y el del almacenamiento se cancelan: en forma no conservativa
el pozo aparece solo como $\sumi{Q\Ceps\pA}$ (consecuencia de la divergencia, no de un $W_-$) y
no hay $\partial_th$.

\subsection{Variación respecto de $h$: adjunto del flujo}
Se perturba $h\to h+\epsilon\,\delta h$, $\delta h(0)=0$. Entonces $\delta\bu=-\bK\nabla\delta h$ y
$\delta(\partial_th)=\partial_t\delta h$. $h$ entra en $\mathcal{L}^{\rm A}$ por tres vías:
\begin{enumerate}[label=\textbf{(\alph*)},nosep]
\item \emph{Ecuación de flujo.} $\langle S_s\partial_t\delta h,\phA\rangle+a_h(\delta h,\phA)$ (lineal en $h$).
\item \emph{Advección} $\int(\bu\cdot\nabla C)\pA$: $\delta=\int(\delta\bu\cdot\nabla C)\pA=-\int\bK\nabla\delta h\cdot(\pA\nabla C)$.
\item \emph{Dispersión} $\int\bD(\bu)\nabla C\cdot\nabla\pA$: $\delta=\int(\partial_{\bu}\bD[\delta\bu]\nabla C)\cdot\nabla\pA
  =-\int\bK\nabla\delta h\cdot\bd$, con
  $\bd=\partial_xC\,\partial_x\pA\,\nabla_{\bu}D_x+\partial_zC\,\partial_z\pA\,\nabla_{\bu}D_z$
  (derivada de $\bD$ respecto de $\bu$).
\item \emph{Almacenamiento} $-\int S_s\partial_th\,C\,\pA$: $\delta=-S_s\int\partial_t\delta h\;C\,\pA$.
\end{enumerate}
Las vías (b) y (c) se agrupan con $\bb=\pA\nabla C+\bd$ y (d) es la contribución del almacenamiento:
\[
\partial_ha_C[\delta h]=-\int_\Omega\bK\bb\cdot\nabla\delta h\;\alm{-\,S_s\int_\Omega\partial_t\delta h\;C\,\pA}.
\]
Integrando por partes en $t$ ($\delta h(0)=0$, $\pA(T)=0$):
\[
\alm{-S_s\!\int_0^T\!\!\!\int_\Omega\partial_t\delta h\,C\pA
=+S_s\!\int_0^T\!\!\!\int_\Omega\delta h\;\partial_t(C\pA)} .
\]
Se suman (a), (b)--(c) y (d), se integra por partes en $t$ el término $S_s\langle\partial_t\delta h,\phA\rangle$
(que da $-S_s\langle\delta h,\partial_t\phA\rangle$ con $\phA(T)=0$) y en espacio los términos con $\nabla\delta h$.
Anulando $\partial_h\mathcal{L}^{\rm A}$ para toda $\delta h$:
\begin{equation}
\boxed{-S_s\partial_t\phA-\nabla\!\cdot(\bK\nabla\phA)
=-\nabla\!\cdot(\bK\bb)\;\alm{-\,S_s\,\partial_t(C\,\pA)}}
\label{eq:adj_flujo}
\end{equation}
con $\phA(T)=0$ y $\bK\nabla\phA\cdot\bn=\bK\bb\cdot\bn$ en $\Gamma_{\rm in}\cup\Gamma_{\rm wall}$. En
$\Gamma_{\rm out}$ la condición libre del estado deja a $\phA$ sobredeterminada
($\phA=0$ y $\bK\nabla\phA\cdot\bn=\bK\bb\cdot\bn$); por eso el adjunto se define a nivel discreto
(\S\ref{sec:discreto}). La fuente verde es nueva: acopla el adjunto del flujo con $C\pA$ aun
con $\bb=0$.

\subsection{Gradiente}
Con $h,C$ y los multiplicadores en su valor de estado/adjunto, $\partial_Q\hat J=\partial_Q\mathcal{L}$.
El control aparece en tres sitios: $-\ell_h(\phA;Q,z_p)=+Q\int g_\varepsilon\phA-\int_{\Gamma_{\rm in}}f_{\rm in}\phA$ (derivada respecto de $Q$:
$+\int g_\varepsilon\phA$), $\beta J_{\rm op}$
(derivada $2\beta d_pQ$) y el kernel $g_\varepsilon(z_p)$; no aparece en el transporte (ya no hay sumidero
$\chi(Q)$):
\begin{equation}
\boxed{\ \nabla_Q\hat J(t)=2\beta\,d_p\,Q(t)+\int_\Omega\Ceps\,\phA\,d\bx\ }
\label{eq:grad_Q}
\end{equation}
\begin{equation}
\boxed{\ \frac{\partial\hat J}{\partial z_p}=\gamma\!\int_0^T\!\!\!\int_\Omega\dot g_\varepsilon C^2
+2\alpha\kappa_p(z_p-z_0)-\beta\!\int_0^TQ^2\,dt
+\int_0^T\!\!Q(t)\!\int_\Omega\dot g_\varepsilon\,\phA\,d\bx\,dt\ }
\label{eq:grad_z}
\end{equation}
con $\dot g_\varepsilon=\partial_{z_p}\Ceps$. Con $Q$ constante a trozos,
$\partial\hat J/\partial Q_n=\int_{I_n}\nabla_Q\hat J\,dt$. Respecto del modelo con sumidero
desaparecen los términos en $\chi'(Q)$ y $\chi$; $\phA$ es otro porque resuelve \eqref{eq:adj_flujo}.

%===========================================================
\section{Adjunto del modelo MRMT}
\label{sec:adj_MRMT}
%===========================================================

\subsection{Variación respecto de $C_m$}
Se perturba $C_m\to C_m+\epsilon\eta$ y se repiten los pasos (a)--(c) de ADR (funcional, tiempo y
$a_C$ con $c_0=\sum_r\lambda_r$, que ya contiene el término $\lambda_rC_m$ del intercambio) y se añade una
contribución: la ecuación inmóvil, multiplicada por $\pim^{\,r}$, contiene $-\lambda_rC_m$, cuya derivada
es $-\sum_r\lambda_r(\pim^{\,r},\eta)$. Por tanto el adjunto móvil es el de ADR más ese acoplamiento:
\begin{equation}
-\phi R\,\partial_t\pM-\nabla\!\cdot(\bD\nabla\pM)-\nabla\!\cdot(\bu\,\pM)
+\Bigl(\sum_r\lambda_r\Bigr)\pM\;\alm{-\,S_s\partial_th\,\pM}-\sum_r\lambda_r\,\pim^{\,r}=-2\gamma\,\Ceps\,C_m,
\label{eq:adj_movil}
\end{equation}
con $\pM(T)=0$ y $(\bD\nabla\pM+\bu\pM)\cdot\bn=0$ en $\partial\Omega$. En forma no conservativa
vale la misma cancelación que en ADR.

\subsection{Variación respecto de $C_{\im,r}$}
$J$ no depende de $C_{\im,r}$. La derivada aporta $-\lambda_r(\pM,\delta C_{\im,r})$ del intercambio móvil
y, de la ecuación inmóvil integrada por partes en $t$,
$-(\phi_{\im,r}R\partial_t\pim^{\,r},\delta C_{\im,r})+\lambda_r(\pim^{\,r},\delta C_{\im,r})$:
\begin{equation}
\boxed{\ \phi_{\im,r}R\,\partial_t\pim^{\,r}=\lambda_r\,(\pim^{\,r}-\pM),
\qquad\pim^{\,r}(\cdot,T)=0\ }
\label{eq:adj_inmovil}
\end{equation}
Ni el sumidero ni el almacenamiento actúan sobre las regiones inmóviles. Con
$\alpha_r=\lambda_r/(\phi_{\im,r}R)$,
\begin{equation}
\pim^{\,r}(t)=\alpha_r\int_t^Te^{-\alpha_r(s-t)}\,\pM(s)\,ds ,
\label{eq:psi_im_memoria}
\end{equation}
una memoria hacia el futuro; el acoplamiento móvil es $\sum_r\lambda_r(\pM-\pim^{\,r})$.

\subsection{Sistema adjunto del transporte MRMT}
\begin{equation}
\left\{\begin{aligned}
&-\phi R\partial_t\pM-\nabla\!\cdot(\bD\nabla\pM)-\bu\cdot\nabla\pM
+\Bigl(\textstyle\sum_r\lambda_r\Bigr)\pM+\sumi{Q\Ceps\pM}-\textstyle\sum_r\lambda_r\pim^{\,r}=-2\gamma\Ceps C_m,\\
&\phi_{\im,r}R\,\partial_t\pim^{\,r}=\lambda_r(\pim^{\,r}-\pM),\quad r=1,\dots,N_r,\\
&\pM(T)=\pim^{\,r}(T)=0 .
\end{aligned}\right.
\label{eq:sistema_adjunto_MRMT}
\end{equation}

\subsection{Adjunto del flujo y gradiente}
$h$ entra por $\bu$, por la ecuación de flujo y por $\alm{-S_s\partial_th\,C_m}$; las regiones inmóviles no
dependen de $h$. Se repite la derivación de ADR con $C=C_m$, $\pA\to\pM$:
\begin{equation}
-S_s\partial_t\phM-\nabla\!\cdot(\bK\nabla\phM)=-\nabla\!\cdot(\bK\bb)\;\alm{-\,S_s\,\partial_t(C_m\,\pM)},
\qquad\bb=\pM\nabla C_m+\bd,\quad\phM(T)=0,
\label{eq:adj_flujo_MRMT}
\end{equation}
y el gradiente es \eqref{eq:grad_Q}--\eqref{eq:grad_z} con $\phA\to\phM$ y $C\to C_m$.

%===========================================================
\section{Adjunto discreto por variante}
\label{sec:discreto}
%===========================================================

Convención: $\mathcal{L}_h=J_h+\sum(\psi^{n+1})^\top(\text{residuo})$,
$J_h=\sum_{n=1}^{N_t}j_n(\mathbf{C}^n)$, $j_n=\gamma\Delta t(\mathbf{C}^n)^\top G\mathbf{C}^n$,
$\nabla j_n=2\gamma\Delta t\,G\mathbf{C}^n$. Los operadores $\mathcal{P},\mathcal{R}$
(y $\mathcal{P}_A,\mathcal{R}_A$) son los de la discretización trapezoidal del transporte sin el
sumidero y con la matriz diagonal de almacenamiento
\[
\alm{\Sigma^n=\frac{S_s}{\Delta t}\,\mathrm{diag}(\mathbf{h}^{n+1}-\mathbf{h}^n)},
\]
que entra como $\alm{-\Sigma^n}$ en la reacción (al ser diagonal, es su propia transpuesta). En FEniCSx la
condición de frontera del adjunto discreto es la natural (transpuesta exacta del estado); BFR impone hoy
$\psi_C=0$ en $\Gamma_{\rm in}\cup\Gamma_{\rm out}$, que no es la transpuesta exacta ni la condición continua
\eqref{eq:adj_fuerte}.

\subsection{ADR ($N_r=0$)}
El paso $n\to n+1$ del estado es $\mathcal{P}^{n+1}\mathbf{C}^{n+1}=\mathcal{R}^{n}\mathbf{C}^n+\mathbf{F}^n$, con
residuo $\mathbf{r}^{n+1}=\mathcal{P}^{n+1}\mathbf{C}^{n+1}-\mathcal{R}^{n}\mathbf{C}^n-\mathbf{F}^n$ y
multiplicador $\boldsymbol\psi_C^{n+1}$. Entonces
$\mathcal{L}_h=J_h+\sum_{m=0}^{N_t-1}(\boldsymbol\psi_C^{m+1})^\top\mathbf{r}^{m+1}$, y $\mathbf{C}^n$ ($1\le n\le N_t$)
aparece en dos residuos: en el del paso que llega a $n$ (con $\mathcal{P}^{n}$) y en el del paso que sale de $n$
(con $-\mathcal{R}^{n}$). Derivando respecto de $\mathbf{C}^n$ e igualando a cero:
\[
\nabla j_n(\mathbf{C}^n)+(\mathcal{P}^{n})^\top\boldsymbol\psi_C^{n}-(\mathcal{R}^{n})^\top\boldsymbol\psi_C^{n+1}=\mathbf{0},
\]
que despejado es \eqref{eq:disc_adr}. Se resuelve hacia atrás desde $n=N_t$ con
$\boldsymbol\psi_C^{N_t+1}=\mathbf{0}$ (no hay paso que salga de $N_t$).
\begin{equation}
(\mathcal{P}^{n})^\top\boldsymbol\psi_C^{n}=(\mathcal{R}^{n})^\top\boldsymbol\psi_C^{n+1}
-\nabla j_n(\mathbf{C}^n),\quad n=N_t,\dots,1,\qquad\boldsymbol\psi_C^{N_t+1}=\mathbf{0}.
\label{eq:disc_adr}
\end{equation}

\subsection{MRMT por bloque (complemento de Schur)}
Estado por paso: $\mathcal{P}^{n+1}\mathbf{C}^{n+1}=\mathcal{R}^{n}\mathbf{C}^n
+\sum_r\mathcal{Q}_r\mathbf{C}_{\im,r}^n+\mathbf{F}^n$ y
$\mathbf{C}_{\im,r}^{n+1}=\kappa_r(\mathbf{C}^{n+1}+\mathbf{C}^n)+(1-2\kappa_r)\mathbf{C}_{\im,r}^n$,
con $\kappa_r=\Delta t\,\omega_r/\eta_r$, $\mathcal{Q}_r=4\phi R\,\omega_r I$. Se añaden a los residuos de $\mathbf{C}$ (multiplicador $\boldsymbol\psi_C^{n+1}$)
los de las regiones inmóviles (multiplicadores $\boldsymbol\psi_r^{n+1}$). Derivando:
\emph{(i)} respecto de $\mathbf{C}_{\im,r}^n$, que aparece en el residuo de $\mathbf{C}$ del paso que sale de $n$
(con $-\mathcal{Q}_r$) y en los residuos inmóviles de los pasos que salen y llegan a $n$
(con $-(1-2\kappa_r)$ y $+1$), se obtiene $\boldsymbol\psi_r^{n}-(1-2\kappa_r)\boldsymbol\psi_r^{n+1}-\mathcal{Q}_r\boldsymbol\psi_C^{n+1}=\mathbf{0}$;
\emph{(ii)} respecto de $\mathbf{C}^n$, que además entra en las dos ecuaciones inmóviles con coeficiente $-\kappa_r$,
se agrega $-\sum_r\kappa_r(\boldsymbol\psi_r^{n}+\boldsymbol\psi_r^{n+1})$ a la ecuación de ADR. Resultado:
\begin{equation}
\left\{\begin{aligned}
\boldsymbol\psi_r^{n}&=(1-2\kappa_r)\,\boldsymbol\psi_r^{n+1}+\mathcal{Q}_r\boldsymbol\psi_C^{n+1},
\qquad \boldsymbol\psi_r^{N_t}=\mathbf{0},\\
(\mathcal{P}^{n})^\top\boldsymbol\psi_C^{n}&=(\mathcal{R}^{n})^\top\boldsymbol\psi_C^{n+1}
+\sum_r\kappa_r\bigl(\boldsymbol\psi_r^{n}+\boldsymbol\psi_r^{n+1}\bigr)-\nabla j_n(\mathbf{C}^n).
\end{aligned}\right.
\label{eq:disc_block}
\end{equation}
Cada paso hacia atrás actualiza primero los $\boldsymbol\psi_r^n$ (diagonales) y luego resuelve un sistema
con $(\mathcal{P}^n)^\top$.

\subsection{MRMT semi-analítico}
Con $e_r=e^{-\alpha_r\Delta t}$, $a_r=\eta_r(1-e_r)$ y $\tilde{\mathbf{C}}_0=\mathbf{C}^*$:
$\mathcal{P}_A^{n+1}\mathbf{C}^*=\mathcal{R}_A^{n}\mathbf{C}^n+\mathbf{F}^n$,
$\tilde{\mathbf{C}}_r=(1-a_r)\tilde{\mathbf{C}}_{r-1}+a_r\mathbf{C}^n_{\im,r}$,
$\mathbf{C}^{n+1}_{\im,r}=(1-e_r)\tilde{\mathbf{C}}_{r-1}+e_r\mathbf{C}^n_{\im,r}$,
$\mathbf{C}^{n+1}=\tilde{\mathbf{C}}_{N_r}$. Con multiplicadores $\boldsymbol\mu_r^{n+1}$ para las
ecuaciones de $\tilde{\mathbf{C}}_r$:
\begin{equation}
\left\{\begin{aligned}
\boldsymbol\mu_{N_r}^{n+1}&=(\mathcal{R}_A^{n+1})^\top\boldsymbol\psi_C^{n+2}-\nabla j_{n+1}(\mathbf{C}^{n+1}),\\
\boldsymbol\mu_{r-1}^{n+1}&=(1-a_r)\,\boldsymbol\mu_r^{n+1}+(1-e_r)\,\boldsymbol\psi_r^{n+1},\qquad r=N_r,\dots,1,\\
(\mathcal{P}_A^{n+1})^\top\boldsymbol\psi_C^{n+1}&=\boldsymbol\mu_0^{n+1},\\
\boldsymbol\psi_r^{n}&=e_r\,\boldsymbol\psi_r^{n+1}+a_r\,\boldsymbol\mu_r^{n+1},
\end{aligned}\right.
\label{eq:disc_semi}
\end{equation}
con $\boldsymbol\psi_C^{N_t+1}=\boldsymbol\psi_r^{N_t}=\mathbf{0}$. Para $N_r=0$ se recupera \eqref{eq:disc_adr}.

\subsection{Adjunto discreto del flujo (las tres variantes)}
Paso de flujo $A_h\mathbf{h}^{n+1}=B_h\mathbf{h}^n+\mathbf{f}^{n+1}$, $\mathbf{f}=-Q\mathbf{g}$. La
fuente del adjunto del flujo tiene dos partes:
\begin{multline}
\boldsymbol\sigma^{m}=\underbrace{-\Bigl(\partial_{\mathbf{h}}\bigl[\mathcal{L}(\bu(\mathbf{h}^m))\mathbf{C}^{*,m}\bigr]\Bigr)^{\!\top}\boldsymbol\psi_C^{m}
-\Bigl(\partial_{\mathbf{h}}\bigl[\mathcal{L}(\bu(\mathbf{h}^m))\mathbf{C}^{m}\bigr]\Bigr)^{\!\top}\boldsymbol\psi_C^{m+1}}_{\text{velocidad}}\\
\alm{+\,\frac{S_s}{\Delta t}\Bigl[\boldsymbol\psi_C^{m+1}\!\circ(\mathbf{C}^{m+1}+\mathbf{C}^m)-\boldsymbol\psi_C^{m}\!\circ(\mathbf{C}^{m}+\mathbf{C}^{m-1})\Bigr]}
\end{multline}
($\mathbf{C}^{*,m}=\mathbf{C}^m$ en ADR y bloque; en el semi-analítico, el par $(\mathbf{C}^*,\mathbf{C}^n)$
del paso A). La parte verde es la derivada temporal discreta de $C\psi_C$, es decir la
discretización de la fuente $-S_s\partial_t(C\psi_C)$ de \eqref{eq:adj_flujo}. Entonces
\begin{equation}
A_h^\top\boldsymbol\psi_h^{m}=B_h^\top\boldsymbol\psi_h^{m+1}-\boldsymbol\sigma^{m},
\quad\boldsymbol\psi_h^{N_t+1}=\mathbf{0},
\qquad
\frac{\partial\hat J}{\partial Q_n}=2\beta d_p\Delta t\,Q_n+\mathbf{g}^\top\boldsymbol\psi_h^{n+1}
\ (\text{con pesos }\theta).
\label{eq:disc_flujo}
\end{equation}
$\partial\hat J/\partial z_p$ es \eqref{eq:grad_z} con las integrales en $t$ sustituidas por sumas con el mismo
$\Delta t$.

\begin{remark}[Convención de signo del código]
El código resuelve $\mathcal{L}=J-\psi\cdot R$ con $R$ en forma de tasa, de modo que
$\psi_{\rm c\acute{o}digo}=-\Delta t\,\psi$ del documento: el $\psi_C$ del código es positivo donde
$C>0$ y el de este documento negativo. Las ecuaciones son las mismas.
\end{remark}

%===========================================================
\section{Algoritmos de gradiente conjugado}
\label{sec:algoritmos}
%===========================================================

Variables de decisión $\mathbf{x}=(Q_1,\dots,Q_{N_t},z_p)$. Cada variante (ADR, MRMT bloque,
MRMT semi) solo cambia el cálculo de $(\hat J,\nabla\hat J)$; el algoritmo exterior es el mismo.

\begin{algorithm}[h]
\caption{Gradiente conjugado no lineal (Polak--Ribière$^+$, Armijo)}
\label{alg:cg}
\begin{algorithmic}[1]
\Require $\mathbf{x}_0$, tolerancia $\varepsilon$, $k_{\max}$, $c=10^{-4}$, $\tau=\tfrac12$, $m_{\max}=20$
\State $(\hat J_0,\mathbf{g}_0)\gets\textsc{EvalJG}(\mathbf{x}_0)$;\quad $\mathbf{d}_0\gets-\mathbf{g}_0$
\For{$k=0,1,\dots,k_{\max}-1$}
  \If{$\|\mathbf{g}_k\|<\varepsilon$} \State \textbf{salir} \EndIf
  \State $\alpha\gets1$
  \For{$i=1,\dots,m_{\max}$} \Comment{búsqueda lineal de Armijo}
    \If{$\hat J(\mathbf{x}_k+\alpha\mathbf{d}_k)\le\hat J_k+c\,\alpha\,\mathbf{g}_k^\top\mathbf{d}_k$} \State \textbf{salir} \EndIf
    \State $\alpha\gets\tau\alpha$
  \EndFor
  \State $\mathbf{x}_{k+1}\gets\Pi\bigl(\mathbf{x}_k+\alpha\mathbf{d}_k\bigr)$ \Comment{$\Pi$ proyección sobre cotas ($^*$)}
  \State $(\hat J_{k+1},\mathbf{g}_{k+1})\gets\textsc{EvalJG}(\mathbf{x}_{k+1})$
  \State $\mathbf{y}_k\gets\mathbf{g}_{k+1}-\mathbf{g}_k$;\quad
         $\beta_k\gets\max\bigl\{0,\ \mathbf{g}_{k+1}^\top\mathbf{y}_k/(\mathbf{g}_k^\top\mathbf{g}_k)\bigr\}$
  \State $\mathbf{d}_{k+1}\gets-\mathbf{g}_{k+1}+\beta_k\mathbf{d}_k$
  \If{$\mathbf{g}_{k+1}^\top\mathbf{d}_{k+1}\ge0$} \State $\mathbf{d}_{k+1}\gets-\mathbf{g}_{k+1}$ \Comment{reinicio ($^*$)} \EndIf
\EndFor
\State \Return $\mathbf{x}_k$
\end{algorithmic}
\end{algorithm}

\begin{remark}
Los pasos marcados con $^*$ (proyección sobre $0\le Q_n\le Q^{\rm sup}_n$, $z_p\in[z_{\min},z_{\max}]$ y reinicio si
$\mathbf{d}$ no es de descenso) \emph{no están} en \texttt{nonlinear\_cg} (\texttt{view/animation.py}), que implementa
Polak--Ribière$^+$ con Armijo sin cotas. La optimización con cotas y con la restricción de suministro se hace en el
código con \texttt{trust-constr} (\texttt{optimize\_bfr}, \texttt{optimize\_fenicsx}), a la que se le entrega el mismo
gradiente adjunto.
\end{remark}

\begin{algorithm}[h]
\caption{\textsc{EvalJG} para ADR}
\label{alg:adr}
\begin{algorithmic}[1]
\Require $\mathbf{x}=(Q_1,\dots,Q_{N_t},z_p)$
\State \textbf{Directo.} $\mathbf{h}^0,\mathbf{C}^0$ dados
\For{$n=0,\dots,N_t-1$}
  \State resolver $A_h\mathbf{h}^{n+1}=B_h\mathbf{h}^n-Q_{n+1}\mathbf{g}$
  \State ensamblar $\mathcal{P}^{n+1},\mathcal{R}^n$ con $\bu(\mathbf{h}^{n+1}),\bu(\mathbf{h}^n)$ y $\Sigma^n$
  \State resolver $\mathcal{P}^{n+1}\mathbf{C}^{n+1}=\mathcal{R}^n\mathbf{C}^n+\mathbf{F}^n$;\ guardar $\mathbf{h}^{n+1},\mathbf{C}^{n+1}$
\EndFor
\State $\hat J\gets\sum_n j_n(\mathbf{C}^n)+\alpha\kappa_p(z_p-z_0)^2+\beta d_p\Delta t\sum_nQ_n^2$
\State \textbf{Adjunto.} $\boldsymbol\psi_C^{N_t+1}\gets\mathbf{0}$,\ $\boldsymbol\psi_h^{N_t+1}\gets\mathbf{0}$
\For{$n=N_t,\dots,1$}
  \State resolver $(\mathcal{P}^{n})^\top\boldsymbol\psi_C^{n}=(\mathcal{R}^{n})^\top\boldsymbol\psi_C^{n+1}-\nabla j_n(\mathbf{C}^n)$ \hfill \eqref{eq:disc_adr}
  \State calcular $\boldsymbol\sigma^{n}$ (velocidad $+$ almacenamiento) y resolver
         $A_h^\top\boldsymbol\psi_h^{n}=B_h^\top\boldsymbol\psi_h^{n+1}-\boldsymbol\sigma^{n}$ \hfill \eqref{eq:disc_flujo}
  \State $g_{Q_n}\gets2\beta d_p\Delta t\,Q_n+\mathbf{g}^\top\boldsymbol\psi_h^{n+1}$
\EndFor
\State $g_{z_p}\gets$ \eqref{eq:grad_z} discretizada
\State \Return $\hat J,\ \mathbf{g}=(g_{Q_1},\dots,g_{Q_{N_t}},g_{z_p})$
\end{algorithmic}
\end{algorithm}

\begin{algorithm}[h]
\caption{\textsc{EvalJG} para MRMT por bloque}
\label{alg:block}
\begin{algorithmic}[1]
\Require $\mathbf{x}$;\ $\mathbf{C}_{\im,r}^0$ dados
\State \textbf{Directo.} como en el Algoritmo~\ref{alg:adr}, pero en cada paso
  $\mathcal{P}^{n+1}\mathbf{C}^{n+1}=\mathcal{R}^n\mathbf{C}^n+\sum_r\mathcal{Q}_r\mathbf{C}_{\im,r}^n+\mathbf{F}^n$
  y luego $\mathbf{C}_{\im,r}^{n+1}=\kappa_r(\mathbf{C}^{n+1}+\mathbf{C}^n)+(1-2\kappa_r)\mathbf{C}_{\im,r}^n$;\ guardar también $\mathbf{C}_{\im,r}^n$
\State \textbf{Adjunto.} $\boldsymbol\psi_C^{N_t+1},\boldsymbol\psi_h^{N_t+1},\boldsymbol\psi_r^{N_t}\gets\mathbf{0}$
\For{$n=N_t,\dots,1$}
  \For{$r=1,\dots,N_r$}
    \State $\boldsymbol\psi_r^{n}\gets(1-2\kappa_r)\boldsymbol\psi_r^{n+1}+\mathcal{Q}_r\boldsymbol\psi_C^{n+1}$ \hfill \eqref{eq:disc_block}
  \EndFor
  \State resolver $(\mathcal{P}^{n})^\top\boldsymbol\psi_C^{n}=(\mathcal{R}^{n})^\top\boldsymbol\psi_C^{n+1}+\sum_r\kappa_r(\boldsymbol\psi_r^{n}+\boldsymbol\psi_r^{n+1})-\nabla j_n$
  \State $\boldsymbol\sigma^n$ y $\boldsymbol\psi_h^{n}$ como en el Algoritmo~\ref{alg:adr}
  \State $g_{Q_n}\gets2\beta d_p\Delta t\,Q_n+\mathbf{g}^\top\boldsymbol\psi_h^{n+1}$
\EndFor
\State $g_{z_p}\gets$ \eqref{eq:grad_z} discretizada;\ \Return $\hat J,\mathbf{g}$
\end{algorithmic}
\end{algorithm}

\begin{algorithm}[h]
\caption{\textsc{EvalJG} para MRMT semi-analítico}
\label{alg:semi}
\begin{algorithmic}[1]
\Require $\mathbf{x}$;\ $\mathbf{C}_{\im,r}^0$ dados;\ $e_r=e^{-\alpha_r\Delta t}$, $a_r=\eta_r(1-e_r)$
\State \textbf{Directo.} en cada paso: resolver $\mathcal{P}_A^{n+1}\mathbf{C}^*=\mathcal{R}_A^n\mathbf{C}^n+\mathbf{F}^n$;
  $\tilde{\mathbf{C}}_0\gets\mathbf{C}^*$;
  \For{$r=1,\dots,N_r$}
    \State $\tilde{\mathbf{C}}_r\gets(1-a_r)\tilde{\mathbf{C}}_{r-1}+a_r\mathbf{C}_{\im,r}^n$;\quad
           $\mathbf{C}_{\im,r}^{n+1}\gets(1-e_r)\tilde{\mathbf{C}}_{r-1}+e_r\mathbf{C}_{\im,r}^n$
  \EndFor
  \State $\mathbf{C}^{n+1}\gets\tilde{\mathbf{C}}_{N_r}$;\ guardar $\mathbf{C}^*$, $\mathbf{C}^{n+1}$, $\mathbf{C}_{\im,r}^{n+1}$
\State \textbf{Adjunto.} $\boldsymbol\psi_C^{N_t+1},\boldsymbol\psi_h^{N_t+1},\boldsymbol\psi_r^{N_t}\gets\mathbf{0}$
\For{$n=N_t-1,\dots,0$}
  \State $\boldsymbol\mu_{N_r}\gets(\mathcal{R}_A^{n+1})^\top\boldsymbol\psi_C^{n+2}-\nabla j_{n+1}(\mathbf{C}^{n+1})$
  \For{$r=N_r,\dots,1$} \Comment{se recorre el splitting al revés}
    \State $\boldsymbol\mu_{r-1}\gets(1-a_r)\boldsymbol\mu_r+(1-e_r)\boldsymbol\psi_r^{n+1}$;\quad
           $\boldsymbol\psi_r^{n}\gets e_r\boldsymbol\psi_r^{n+1}+a_r\boldsymbol\mu_r$ \hfill \eqref{eq:disc_semi}
  \EndFor
  \State resolver $(\mathcal{P}_A^{n+1})^\top\boldsymbol\psi_C^{n+1}=\boldsymbol\mu_0$
  \State $\boldsymbol\sigma^{n+1}$ con el par $(\mathbf{C}^*,\mathbf{C}^n)$ y resolver $\boldsymbol\psi_h^{n+1}$ como en el Algoritmo~\ref{alg:adr}
  \State $g_{Q_{n+1}}\gets2\beta d_p\Delta t\,Q_{n+1}+\mathbf{g}^\top\boldsymbol\psi_h^{n+1}$
\EndFor
\State $g_{z_p}\gets$ \eqref{eq:grad_z} discretizada;\ \Return $\hat J,\mathbf{g}$
\end{algorithmic}
\end{algorithm}

\paragraph{Costo.} Cada $\textsc{EvalJG}$ cuesta un barrido directo y uno adjunto, es decir dos
resoluciones de flujo y dos de transporte por paso, independientemente del número de controles ($N_t+1$);
MRMT por bloque añade operaciones diagonales por región y el semi-analítico un barrido de $N_r$ pasos
diagonales. Cada iteración de CG usa una evaluación de $\hat J$ por intento de la búsqueda de Armijo y un
$\textsc{EvalJG}$ completo en el punto aceptado.

%===========================================================
\section{Verificaciones ejecutadas}
\label{sec:verif}
%===========================================================

Todo lo siguiente se ejecutó; lo que no figura aquí no se verificó.
\begin{itemize}[nosep]
\item \textbf{Adjunto discreto de FEniCSx frente a diferencias finitas centradas} (solo $\gamma J_m$, la parte
  económica es analítica y se resta; 3 pasos; $C$ uniforme y modulada; con y sin almacenamiento): ADR, MRMT
  semi-analítico y MRMT por bloque coinciden con error relativo $10^{-6}$--$10^{-10}$.
\item \textbf{Ecuaciones continuas frente al código (ADR, FEniCSx, $\theta_C=1$):} la forma débil de
  \eqref{eq:adj_fuerte} escrita aquí, evaluada sobre el $\psi_C$ del código con $\pA=-\psi_{\rm c\acute{o}digo}/\Delta t$,
  da un residuo relativo $\sim10^{-15}$; sin el término de almacenamiento sube a $\sim10^{-3}$ y sin la
  advección a $\sim10^{-2}$. Para \eqref{eq:adj_flujo}, la descomposición de $(\partial_hR^C)^\top\psi_C$ en los
  términos $\int\bK\pA\nabla C\cdot\nabla v$, $\int\bK\bd\cdot\nabla v$ y $\int\phi C\pA v/\Delta t$ da
  coeficientes $-1.000000$ y residuo $5\times10^{-16}$.
\item \textbf{No verificado:} el adjunto de BFR (con $C$ no uniforme no coincide con las diferencias finitas: factor
  2--7), las ecuaciones continuas de MRMT, el balance de masa numérico, la convergencia en malla y tiempo, y una
  optimización completa convergida.
\end{itemize}

\end{document}
